% GENERATED FILE. Edit canonical lessons, then run npm run generate:books.
% canonical-source-hash: fnv1a64:29a82d1e894ac0ad
% canonical-lessons: IT-C462-vetta, IT-C462-pendio, IT-C462-burrone, IT-C462-caverna, IT-C462-arrampicarsi

\chapter{In montagna}
\label{ch:it-in-montagna}
\hlchaptermodality{462}

\begin{chapteropening}
\textbf{By the end of this chapter:} \emph{I can describe a mountain walk: the summit, the slope, a ravine, a cave and climbing.}

Complete the last lesson of chapter 462: I can describe a mountain walk: the summit, the slope, a ravine, a cave and climbing.
\end{chapteropening}

\section[la vetta]{la vetta — the summit, the peak}

\label{lesson:IT-C462-vetta}



\begin{quote}
Before the new one: say the Italian for to prune, then the Italian for the bark.
\end{quote}

\subsection*{You'll want to know: la vetta}
\textbf{la vetta} — \textquotedblleft{}the summit, the peak\textquotedblright{}.

\begin{quote}
Dalla vetta si vedono due laghi.
\end{quote}

\begin{quote}
\emph{From the summit you can see two lakes.}
\end{quote}

Like \emph{la cima}, a little more literary.

\begin{cousinweb}
\emph{vetta} is traced in EtymDB to Latin \emph{vitta} 'headband, ribbon'; it is inherited, and first meant the tip of a branch.
\end{cousinweb}

\begin{grammarlens}[title={What you've built}]
One more word for telling, explaining and describing at B1.
\end{grammarlens}

\subsection*{Your turn}
\begin{itemize}
\raggedright
  \item \emph{Say it:} \emph{la vetta}
  \item \emph{Say it:} \emph{la vetta}, once more
  \item \emph{Say it:} a sentence of your own with \emph{la vetta}
  \item \emph{Recall:} say the Italian for to prune, then the Italian for the bark, then say \emph{la vetta} again
\end{itemize}

\begin{tcolorbox}[breakable,colback=teal!4,colframe=teal!35!black,title={Before you move on}]
What does la vetta mean? (\textquotedblleft{}The summit, the peak\textquotedblright{}.) Say it once more.
\end{tcolorbox}
\section[il pendio]{il pendio — the slope, the hillside}

\label{lesson:IT-C462-pendio}



\begin{quote}
Before the new one: say the Italian for the bark, then the Italian for the summit.
\end{quote}

\subsection*{You'll want to know: il pendio}
\textbf{il pendio} — \textquotedblleft{}the slope, the hillside\textquotedblright{}.

\begin{quote}
Le case sono sul pendio della collina.
\end{quote}

\begin{quote}
\emph{The houses are on the slope of the hill.}
\end{quote}

\emph{In pendio} — sloping.

\begin{cousinweb}
\emph{pendio} is built on \emph{pendere} 'to lean, to slope', which is inherited from Latin \emph{pendēre} 'to hang' (Lewis \& Short).
\end{cousinweb}

\begin{grammarlens}[title={What you've built}]
One more word for telling, explaining and describing at B1.
\end{grammarlens}

\subsection*{Your turn}
\begin{itemize}
\raggedright
  \item \emph{Say it:} \emph{il pendio}
  \item \emph{Say it:} \emph{il pendio}, once more
  \item \emph{Say it:} a sentence of your own with \emph{il pendio}
  \item \emph{Recall:} say the Italian for the bark, then the Italian for the summit, then say \emph{il pendio} again
\end{itemize}

\begin{tcolorbox}[breakable,colback=teal!4,colframe=teal!35!black,title={Before you move on}]
What does il pendio mean? (\textquotedblleft{}The slope, the hillside\textquotedblright{}.) Say it once more.
\end{tcolorbox}
\section[il burrone]{il burrone — the ravine, the gorge}

\label{lesson:IT-C462-burrone}



\begin{quote}
Before the new one: say the Italian for the summit, then the Italian for the slope.
\end{quote}

\subsection*{You'll want to know: il burrone}
\textbf{il burrone} — \textquotedblleft{}the ravine, the gorge\textquotedblright{}.

\begin{quote}
Non avvicinarti al burrone!
\end{quote}

\begin{quote}
\emph{Do not go near the ravine!}
\end{quote}

Deep and steep.

\begin{cousinweb}
\emph{burrone} is built from \emph{borro} 'ditch, gully' and the big-size ending \emph{-one} (EtymDB); the deeper origin of \emph{borro} is not settled.
\end{cousinweb}

\begin{grammarlens}[title={What you've built}]
One more word for telling, explaining and describing at B1.
\end{grammarlens}

\subsection*{Your turn}
\begin{itemize}
\raggedright
  \item \emph{Say it:} \emph{il burrone}
  \item \emph{Say it:} \emph{il burrone}, once more
  \item \emph{Say it:} a sentence of your own with \emph{il burrone}
  \item \emph{Recall:} say the Italian for the summit, then the Italian for the slope, then say \emph{il burrone} again
\end{itemize}

\begin{tcolorbox}[breakable,colback=teal!4,colframe=teal!35!black,title={Before you move on}]
What does il burrone mean? (\textquotedblleft{}The ravine, the gorge\textquotedblright{}.) Say it once more.
\end{tcolorbox}
\section[la caverna]{la caverna — the cave, the cavern}

\label{lesson:IT-C462-caverna}



\begin{quote}
Before the new one: say the Italian for the slope, then the Italian for the ravine.
\end{quote}

\subsection*{You'll want to know: la caverna}
\textbf{la caverna} — \textquotedblleft{}the cave, the cavern\textquotedblright{}.

\begin{quote}
Nella caverna fa freddo anche d'estate.
\end{quote}

\begin{quote}
\emph{In the cave it is cold even in summer.}
\end{quote}

Bigger and deeper than \emph{una grotta}.

\begin{cousinweb}
\emph{caverna} is a learned borrowing from Latin \emph{caverna} 'cave, hollow' (EtymDB), from \emph{cavus} 'hollow'.
\end{cousinweb}

\begin{grammarlens}[title={What you've built}]
One more word for telling, explaining and describing at B1.
\end{grammarlens}

\subsection*{Your turn}
\begin{itemize}
\raggedright
  \item \emph{Say it:} \emph{la caverna}
  \item \emph{Say it:} \emph{la caverna}, once more
  \item \emph{Say it:} a sentence of your own with \emph{la caverna}
  \item \emph{Recall:} say the Italian for the slope, then the Italian for the ravine, then say \emph{la caverna} again
\end{itemize}

\begin{tcolorbox}[breakable,colback=teal!4,colframe=teal!35!black,title={Before you move on}]
What does la caverna mean? (\textquotedblleft{}The cave, the cavern\textquotedblright{}.) Say it once more.
\end{tcolorbox}
\section[arrampicarsi]{arrampicarsi — to climb (using hands and feet)}

\label{lesson:IT-C462-arrampicarsi}



\begin{quote}
Before the new one: say the Italian for the ravine, then the Italian for the cave.
\end{quote}

\subsection*{You'll want to know: arrampicarsi}
\textbf{arrampicarsi} — \textquotedblleft{}to climb (using hands and feet)\textquotedblright{}.

\begin{quote}
I ragazzi si arrampicano sulle rocce.
\end{quote}

\begin{quote}
\emph{The boys climb on the rocks.}
\end{quote}

\emph{L'arrampicata} — rock climbing.

\begin{cousinweb}
\emph{arrampicarsi} is built on \emph{rampa} 'claw, paw', which is usually traced to a Germanic word that came into Italian as a borrowing: climbing with your claws.
\end{cousinweb}

\begin{grammarlens}[title={What you've built}]
One more word for telling, explaining and describing at B1.
\end{grammarlens}

\subsection*{Your turn}
\begin{itemize}
\raggedright
  \item \emph{Say it:} \emph{arrampicarsi}
  \item \emph{Say it:} \emph{arrampicarsi}, once more
  \item \emph{Say it:} a sentence of your own with \emph{arrampicarsi}
  \item \emph{Recall:} say the Italian for the ravine, then the Italian for the cave, then say \emph{arrampicarsi} again
\end{itemize}

\begin{tcolorbox}[breakable,colback=teal!4,colframe=teal!35!black,title={Before you move on}]
What does arrampicarsi mean? (\textquotedblleft{}To climb (using hands and feet)\textquotedblright{}.) Say it once more.
\end{tcolorbox}
